\section*{Introducción a la probabilidad}

La \emph{probabilidad} es la rama de las matemáticas que estudia los fenómenos aleatorios, es decir, aquellos cuyos resultados no pueden determinarse de antemano con certeza. A diferencia de los fenómenos deterministas (como la caída de un objeto bajo la gravedad), los fenómenos aleatorios presentan variabilidad inherente que solo puede describirse mediante modelos probabilísticos.

\subsection*{¿Por qué estudiar probabilidad?}

La probabilidad es fundamental en múltiples áreas del conocimiento:

\begin{itemize}
	\item \textbf{Ciencias naturales:} Modelado de fenómenos cuánticos, genética, evolución, termodinámica estadística.
	\item \textbf{Ingeniería:} Confiabilidad de sistemas, teoría de colas, procesamiento de señales.
	\item \textbf{Economía y finanzas:} Valoración de opciones, gestión de riesgos, modelos de mercado.
	\item \textbf{Ciencias de la computación:} Algoritmos aleatorizados, aprendizaje automático, criptografía.
	\item \textbf{Medicina:} Ensayos clínicos, epidemiología, diagnóstico probabilístico.
	\item \textbf{Vida cotidiana:} Toma de decisiones bajo incertidumbre, juegos de azar, predicciones meteorológicas.
\end{itemize}

\subsection*{Relación con la estadística}

La probabilidad y la estadística están íntimamente relacionadas:

\begin{itemize}
	\item La \textbf{probabilidad} proporciona el marco teórico para modelar fenómenos aleatorios y calcular la likelihood de eventos.
	\item La \textbf{estadística inferencial} utiliza datos observados para hacer inferencias sobre poblaciones, utilizando los modelos probabilísticos como fundamento.
\end{itemize}

En otras palabras, la probabilidad nos dice qué tan probable es observar ciertos datos dado un modelo, mientras que la estadística nos ayuda a inferir el modelo a partir de los datos observados.

\subsection*{Conceptos fundamentales}

En este capítulo desarrollaremos los siguientes conceptos:

\begin{enumerate}
	\item \textbf{Experimentos aleatorios y espacio muestral:} Definiremos qué es un experimento aleatorio y cómo describir todos sus posibles resultados.
	\item \textbf{Eventos y operaciones:} Estudiaremos cómo combinar eventos mediante unión, intersección y complemento.
	\item \textbf{Axiomas de probabilidad:} Presentaremos el enfoque axiomático moderno de Kolmogorov.
	\item \textbf{Probabilidad condicional e independencia:} Analizaremos cómo la ocurrencia de un evento afecta la probabilidad de otro.
	\item \textbf{Teorema de Bayes:} Una herramienta fundamental para actualizar probabilidades con nueva información.
	\item \textbf{Variables aleatorias:} Funciones que asignan valores numéricos a los resultados de experimentos aleatorios.
	\item \textbf{Distribuciones especiales:} Modelos probabilísticos comúnmente utilizados (binomial, normal, Poisson).
\end{enumerate}

\subsection*{Un poco de historia}

El estudio matemático de la probabilidad comenzó en el siglo XVII con correspondencia entre Blaise Pascal y Pierre de Fermat sobre problemas de juegos de azar. Posteriormente, Jacob Bernoulli demostró la Ley de los Grandes Números, y en el siglo XX, Andrey Kolmogorov estableció los axiomas modernos de la probabilidad en 1933, proporcionando una base matemática rigurosa.

\begin{observacion}
	A lo largo de este capítulo, utilizaremos ejemplos concretos (lanzamiento de monedas, dados, extracción de cartas) para ilustrar los conceptos. Sin embargo, estos mismos principios se aplican a problemas mucho más complejos en ciencia, ingeniería y toma de decisiones.
\end{observacion}
